% Appendix (IEEE: before the acknowledgment), then the back matter.
\appendices
\section{Additional Tables}
\label{app:extra}
The secondary fidelity metrics (band-limited SSIM, MS-SSIM, LPIPS, raw metrics and the high-frequency energy ratio) are in
Table~\ref{tab:szch_heldout}; the per-weight and per-encoder results in Tables~\ref{tab:dose_response} and~\ref{tab:encoders}. The
deviations from the analysis plans, further results of Studies~II and~III and the per-seed results are in the supplementary material.

\begin{table*}[!t]
\centering
\caption{Single-seed results of the final models. The ensembles average seeds 0 to 2; SwinV2-T seed 3 is the reserve.}
\label{tab:app_seeds}
\restable[\textwidth]{app_seeds}
\end{table*}

\begin{table*}[!t]
\centering
\caption{Exact-truth audit on the Qure-projection test split (430 projections; lesions added to the soft component): absorption, lesion
contrast ratio of the soft estimate and bone MAE. The SwinV2-T twin row averages the four seeds of the dose-response arm ($K=4$); the other
trained models use seeds 0--2.}
\label{tab:app_qp}
\restable[\textwidth]{app_qp_audit}
\end{table*}

\begin{table}[!t]
\centering
\caption{Audit amplitude sweep on the 450 SZCH sites (development data, models of 27 October 2026).}
\label{tab:app_amplitude}
\restable{app_amplitude}
\end{table}

\begin{table}[!t]
\centering
\caption{Encoder screen at \num{100000} samples, twin weight 16 (development data, rule R1).}
\label{tab:app_screen}
\restable{app_encoder_screen}
\end{table}

\clearpage   % place all result and appendix floats before the back matter and the reference list
\section*{Acknowledgment}
We thank the providers of the data used here: Qure.ai for the CT projections and the released bone-suppression weights, the CT-RATE team,
The Cancer Imaging Archive and the LIDC-IDRI investigators, the authors of SZCH-X-Rays, VinDr (VinDr-RibCXR and VinDr-CXR), the Japanese
Society of Radiological Technology (JSRT), the Node21, TBX11K and NIH ChestX-ray14 teams, and Dr.~M~R~Khan Shishu Hospital \& Institute of
Child Health, Dhaka, for BDCXR-3257 (ethical approval REC-FSIT-2025/No:~11928). Training ran on Kaggle GPU quota. \todo{funding and
colleagues}

\section*{Ethics Statement}
This is a retrospective study of existing images. Public data sets were used under their licences, VinDr-RibCXR under its data use
agreement, and BDCXR-3257 under the ethical approval REC-FSIT-2025/No:~11928 of Daffodil International University and a data use agreement
with the providing hospital. No new patient data were collected, and no image of a licensed or restricted set is shown.

\section*{Data and Code Availability}
The public data sets are available from their providers (Table~\ref{tab:data}); CT-RATE \cite{hamamci2026ctrate} is available from its
providers under CC BY-NC-SA~4.0, and no data derived from it beyond metrics and averaged heat maps are shared. VinDr-RibCXR was used under
its data use agreement and is not redistributed. BDCXR-3257 is public on figshare under CC BY 4.0 \cite{durjoy2026bdcxr}; the authors
obtained it under a data use agreement with the providing hospital and follow its terms, so no image of either set is shown and nothing
derived from it is redistributed. VinBigData was processed only in cloud accounts that accepted the competition rules. Only metrics left
these environments, and no image, mask, map or model derived from VinDr-RibCXR, VinBigData or BDCXR-3257 images is released.

\emph{Analysis plans.} The analysis plans (SAP.md and SAP\_v2.md), the split files, the weights and the code commit were hashed into
\path{freeze.json} and \path{freeze_v2.json} before any held-out or locked evaluation: \path{freeze.json} on 27 October 2026 (code commit
\texttt{8e9e3be}) and \path{freeze_v2.json} on 1 November 2026 (commit \texttt{dd52b6a}); the BDCXR-3257 addendum was committed on
3 November 2026 before any model saw those images, and the classifier study's plan (SAP\_v3.md), code and split files were hashed into
\path{freeze_v3.json} on 3 November 2026 (commit \texttt{c39aad5}) before its encoder screen. The cross-modality study's plan (SAP\_v4.md),
its design, code and CT subset list were hashed into \path{freeze_v4.json} on 8 November 2026 before its first training run. The
exploratory site calibration's plan (SAP\_v5.md), code (\path{s4/site_cal.py}, \path{tools/szch_site_cal.py}) and coefficients fitted on the
24 development pairs (\path{results/site_cal_coef.json}) were hashed into \path{freeze_v5.json} on 17 November 2026 (commit \texttt{e43c59a})
before any held-out pair was scored with them. The plans were kept in the project's version control, not deposited in a public registry;
they and the hash files are supplementary material.

\emph{Reproducing the numbers.} Every table and figure is produced by a script from saved result files; no number in the text is typed by
hand. From the saved results, \texttt{python} \path{tools/final_stats.py} writes \path{results/final/stats_held.json} and the main tables,
\texttt{python} \path{tools/paper_extras.py} \texttt{--audit-map} the ablation table and the dose-response, audit-map, downstream and
reliability figures, \texttt{python} \path{tools/fig_pareto.py} \texttt{final} and \path{tools/fig_real_qualitative.py} the remaining
figures, and \texttt{python} \path{tools/bdcxr_stats.py} the BDCXR-3257 tables; the evaluations themselves are run by
\path{tools/final_local.sh} (locally or in cloud kernels). For the classifier study, \texttt{python} \path{tools/hybrid_stats.py} writes
\path{results/final_hybrid_stats.json} and its tables, and \texttt{python} \path{tools/generate_all_thesis_figures.py} its figures; its 175
training runs used the Kaggle kernels of the repository. For the cross-modality study, \texttt{python} \path{tools/xm_stats.py} writes
\path{results/final_xm_stats.json}, \path{tools/xm_tables.py} its tables and \path{tools/fig_v4.py} its figures. For the exploratory site
calibration, \texttt{python} \path{tools/szch_site_cal.py} with the steps \texttt{heldout}, \texttt{audit}, \texttt{audit\_stats} and
\texttt{rib} runs it, and \path{tools/site_cal_tables.py} writes its tables. All of this is in the repository
\url{https://github.com/HTanvir0011/Thesis-full-project} (release v1.0.0).

\emph{Compute.} Training used Kaggle T4 GPUs: a run of \num{300000} samples took 6 to 7 GPU-hours on one T4 (two runs share a two-GPU
session), a screening run of \num{100000} samples about 2 GPU-hours; the per-run times are in the training logs. The locked-set evaluation
ran on a consumer GPU and on cloud T4 GPUs with identical code; the audit sites of every model were pinned to the same lung masks, and model
outputs differ between the two GPU types by at most 0.0024\,dB.

\emph{Release.} The code is released under the MIT licence, and the three checkpoints of the ensemble and a model card for research use
under CC BY-NC-SA 4.0, the licence of the training projections, as release v1.0.0 of the repository, whose SHA-256 checksums equal the hashes
frozen before the evaluation. The reporting follows the
CLAIM 2024 checklist \cite{tejani2024claim}, which is submitted with the manuscript.

\section*{CRediT Authorship Contribution Statement}
\todo{author contributions in CRediT terms, for example conceptualisation, methodology, software, formal analysis, writing}

\section*{Declaration of Competing Interest}
\todo{competing interests, or: The authors declare that they have no known competing financial interests or personal relationships that
could have appeared to influence the work reported in this paper.}

\section*{Declaration of Generative AI Use}
During the preparation of this work the authors used an AI assistant (Claude, Anthropic) for code, analysis automation and drafting of the
text. After using this tool, the authors reviewed and edited the content as needed and take full responsibility for the content of the
publication. \todo{confirm or adapt this statement to the actual use before submission}
